\subsection{Sparse edges in growing dimension}
\label{sec:proof-program}

This subsection extends the spatial exclusion required by
Theorem~\ref{thm:main-support-law} from fixed to growing dimension below the
simplex regime. We split the sample into a sparse upper radial cloud and a
truncated lower cloud. Angular separation controls the upper cloud, while
dimension-uniform concentration controls the lower cloud.

Fix a sequence \(A_n\to\infty\) and define the upper radial cloud
\begin{equation}
\label{eq:upper-cloud}
  \mathcal I_n(A_n)
  =
  \{i:\xi_{n,i}\ge-A_n\}.
\end{equation}
The expected size of this set is \(e^{A_n+o(A_n)}\) whenever
\(A_n=o(L_n)\). Write \(B_n=B_n^{\rm hi}+B_n^{\rm lo}\) for the two sums
in \eqref{eq:bump-representation}, split according to
\(\mathcal I_n(A_n)\).

\subsubsection{Separation of the upper radial cloud}

Conditional on their radii, the uncentered Gaussian directions are
independent and uniform on \(S^{d_n-1}\). If
\(\#\mathcal I_n(A_n)=\exp\{O_{\Pbb}(A_n)\}\) and \(A_n=o(d_n)\), spherical
concentration gives
\begin{equation}
\label{eq:direction-separation}
  \max_{i\ne j\in\mathcal I_n(A_n)}
  \left|
    \frac{q_i}{|q_i|}\cdot\frac{q_j}{|q_j|}
  \right|
  =o_{\Pbb}(1).
\end{equation}
Since the upper centers have radius
\(\sqrt{2L_n}\{1+o_{\Pbb}(1)\}\),
\eqref{eq:direction-separation} implies
\[
  \min_{i\ne j\in\mathcal I_n(A_n)}|q_i-q_j|^2
  =
  4L_n\{1+o_{\Pbb}(1)\}.
\]
Thus no location \(w\) can lie within \(o(\sqrt{L_n})\) of two upper
centers. The contribution of all nonnearest upper bumps is exponentially
small in \(L_n\).

Centering by \(\bar Z_n\) perturbs every direction by
\(O_{\Pbb}(\sqrt{d_n/n})\) in Euclidean norm. Under
\eqref{eq:centering-condition} and the rate ranges considered below, this
does not affect \eqref{eq:direction-separation}.

\subsubsection{The lower cloud and the inner field}

At the critical normalization, every low bump has height at most
\begin{equation}
\label{eq:low-jump}
  \max_{i\notin\mathcal I_n(A_n)}h_i
  \le
  \exp\{-\delta_n(A_n+s_n)\}.
\end{equation}
Before empirical centering, let
\[
  B_n^\circ(w)
  =
  \frac1n\sum_{i=1}^n
  \exp\left\{
    w\cdot\frac{Z_i}{\sqrt{\gamma_n}}-\frac{|w|^2}{2}
  \right\}.
\]
For a fixed \(w\), its population mean is
\begin{equation}
\label{eq:population-field}
  \E B_n^\circ(w)
  =
  \exp\left\{-\frac{\eps_n|w|^2}{2}\right\}.
\end{equation}
Thus the deterministic background is already negligible near
\(|w|=\sqrt{2L_n}\) whenever \(d_n\to\infty\), because then
\(\eps_nL_n\to\infty\) at the critical scale.

The proof below establishes the following low-cloud estimate.  Its precision
\(o_{\Pbb}(\delta_n)\) is essential in superlogarithmic dimension: for a
deterministic radius \(r_n=o(\sqrt{L_n})\),
\begin{equation}
\label{eq:outer-low-target}
  \sup_{|w|\ge r_n}B_n^{\rm lo}(w)
  =o_{\Pbb}(\delta_n),
\end{equation}
provided \(A_n\) is chosen so that
\begin{equation}
\label{eq:A-conditions}
  A_n=o(L_n\wedge d_n),
  \qquad
  \delta_nA_n-\log(d_nL_n)\longrightarrow\infty.
\end{equation}

For fixed \(w\), \eqref{eq:low-jump}, \eqref{eq:population-field}, and
Bernstein's inequality give much more than \eqref{eq:outer-low-target}. The
work is to transfer that bound over \(w\). A direct Euclidean net of the
relevant ball has logarithmic size \(O(d_n\log L_n)\); the second condition
in \eqref{eq:A-conditions} is designed to make the bounded-jump exponent
dominate this entropy. Radial peeling avoids an unnecessarily fine net
near the origin.

The inner region contains the deterministic equality \(B_n(0)=1\), so it
cannot satisfy \eqref{eq:outer-low-target}.  The corresponding target is
strict concavity away from zero:
\begin{equation}
\label{eq:inner-target}
  \Pbb\left\{
    \sup_{0<|w|\le r_n}B_n(w)<1
  \right\}
  \longrightarrow1.
\end{equation}

At the origin,
\[
  \nabla B_n(0)=0,
  \qquad
  \nabla^2B_n(0)
  =
  \frac1n\sum_i
  \frac{\widetilde Z_i\widetilde Z_i^\top}{\gamma_n}
  -I_{d_n}.
\]
Its expectation is \(-\eps_nI_{d_n}+o(1)\). Standard covariance concentration
therefore gives a negative Hessian whenever
\(\sqrt{d_n/n}=o(\eps_n)\). The remaining task is to propagate this local
curvature through the interval up to \(r_n\) using a clipped exponential
process. This is easier than the outer-edge estimate because the population
gap in \eqref{eq:population-field} grows monotonically with \(|w|\).

Equations~\eqref{eq:outer-low-target} and \eqref{eq:inner-target} give the
field exclusion used later in
Proposition~\ref{prop:buffered-support-below-simplex}. On the event
\(\max_i\xi_{n,i}\le s_n-\eta\), every upper bump has height at most
\[
  e^{-\delta_n\eta},
\]
and hence lies below one by at least \(c_\eta\delta_n\) whenever
\(0<\delta_n\le1\).
At the center of any upper bump, the separation estimate makes all other
upper bumps \(o_{\Pbb}(\delta_n)\), and \eqref{eq:outer-low-target} makes the
low-cloud contribution \(o_{\Pbb}(\delta_n)\). Hence the outer field is below one, while
\eqref{eq:inner-target} handles the inner field.

\subsubsection{The direct truncation range}

The truncation and angular-separation conditions can be imposed
simultaneously in the range
\begin{equation}
\label{eq:direct-truncation-range}
  \log\log n\ll d_n,
  \qquad
  d_n\{\log d_n\}^2=o\{(\log n)^3\},
  \qquad
  d_n\log n=o(n).
\end{equation}
Indeed, in the worst part of this range,
\(\delta_n\asymp\sqrt{L_n/d_n}\). Conditions
\eqref{eq:A-conditions} admit a choice of \(A_n\) precisely when
\[
  \delta_n^{-1}\log(d_nL_n)=o(L_n\wedge d_n),
\]
which is implied by \eqref{eq:direct-truncation-range}. The lower bound
\(d_n\gg\log\log n\) ensures \(A_n=o(d_n)\), so the upper radial points are
uniformly separated. Smaller dimensions are handled below by constant-mesh
estimates. At the upper end, where \(\delta_n\) is too small for
\eqref{eq:A-conditions}, Section~\ref{sec:explicit-regime} replaces spatial
truncation by the entropy dual.

\subsubsection{Slowly growing dimension}
\label{sec:slow-dimension}

We first treat growing dimensions below \((\log L_n)^2\). A constant-mesh
argument is enough; its key input is an exact semiconvexity property of every
Gaussian-bump field. Fixed dimensions use the moving-shell estimates in
Section~\ref{sec:fixed-d-patches}.

\begin{lemma}
\label{lem:slow-truncated-shell}
Let
\[
 F(w)=\sum_{j=1}^m c_j
 \exp\left\{w\cdot q_j-\frac{|w|^2}{2}\right\},
 \qquad c_j>0.
\]
Then
\begin{equation}
\label{eq:semiconvex-log-field}
 \nabla^2\log F(w)=\operatorname{Cov}_{p_w}(q_j)-I\succeq-I,
\end{equation}
where \(p_w\) is proportional to
\(c_j e^{w\cdot q_j-|w|^2/2}\).  Every global maximizer \(w_*\) belongs
to the convex hull of \(q_1,\ldots,q_m\), and
\begin{equation}
\label{eq:semiconvex-net-transfer}
 F(w)\ge F(w_*)e^{-|w-w_*|^2/2}
 \qquad(w\in\R^d).
\end{equation}
\end{lemma}

\begin{proof}
Differentiating \(\log F\) gives \eqref{eq:semiconvex-log-field}.  At a
maximizer, \(0=\nabla\log F(w_*)=\sum_jp_{w_*}(j)q_j-w_*\), which proves
the convex-hull assertion.  Integrating the lower Hessian bound along the
segment from \(w_*\) to \(w\) proves
\eqref{eq:semiconvex-net-transfer}.
\end{proof}

Fix \(\alpha>0\), put
\[
 Q_i=\frac{Z_i}{\sqrt{\gamma_n}},\qquad
 R_{n,\alpha}=\sqrt{2(L_n-\alpha)},
\]
and define the independent truncated field
\[
 S_{n,\alpha}(w)
 =\frac1n\sum_{i=1}^n
 e^{w\cdot Q_i-|w|^2/2}
 \mathbf1_{\{|Q_i|\le R_{n,\alpha}\}}.
\]
The next lemma gives the only outer-field estimate needed below.  It uses
only bounded summands, their population mean, and a constant Euclidean net.

\begin{lemma}
\label{lem:slow-outside-shell}
Suppose
\[
 d_n\longrightarrow\infty,\qquad
 d_n\le(\log L_n)^2,\qquad s_n=O(1).
\]
For every fixed \(\alpha>0\), choose a fixed \(\zeta>0\) with
\(\zeta^2/2<\alpha\), let \(\mathcal G_n\) be a \(\zeta\)-net of
\(B(0,R_{n,\alpha})\), and fix
\[
 0<\rho<e^{-\zeta^2/2}-e^{-\alpha}.
\]
Then
\begin{equation}
\label{eq:slow-outside-shell}
 \Pbb\left\{
  \max_{\substack{w\in\mathcal G_n\\
       |w|\ge R_{n,\alpha}/2-\zeta}}
  S_{n,\alpha}(w)
  \ge e^{-\zeta^2/2}-\rho
 \right\}\longrightarrow0.
\end{equation}
\end{lemma}

\begin{proof}
For fixed \(w\), every summand of \(S_{n,\alpha}(w)\) is at most
\begin{equation}
\label{eq:slow-jump-simple}
 M_n(w)=\exp\left\{-\alpha
 -\frac{(R_{n,\alpha}-|w|)^2}{2}\right\},
\end{equation}
while
\begin{equation}
\label{eq:slow-mean-simple}
 \E S_{n,\alpha}(w)\le e^{-\eps_n|w|^2/2}.
\end{equation}
If independent nonnegative summands have sum of means \(\mu\), are bounded
by \(M\), and \(t>\mu\), the exponential-Chernoff bound is
\begin{equation}
\label{eq:bounded-poisson-chernoff}
 \Pbb\left\{\sum_jX_j\ge t\right\}
 \le
 \exp\left[-\frac1M
 \left\{t\log\frac t\mu-t+\mu\right\}\right].
\end{equation}
Apply this with \(t=e^{-\zeta^2/2}-\rho\).  Uniformly for
\(R_{n,\alpha}/2-\zeta\le|w|\le R_{n,\alpha}\),
\[
 -\log\Pbb\{S_{n,\alpha}(w)\ge t\}
 \ge(1-o(1))e^\alpha t\,\eps_nL_n.
\]
Indeed, the infimum of
\(e^{(R_{n,\alpha}-r)^2/2}r^2\) over this interval is
\((1-o(1))R_{n,\alpha}^2\).

In the present range, \(k_n=d_n/2\to\infty\) and
\begin{equation}
\label{eq:slow-entropy-match}
 \eps_nL_n
 =(1+o(1))k_n\log\frac{L_n}{k_n}
 =(1-o(1))k_n\log L_n.
\end{equation}
On the other hand,
\[
 \log|\mathcal G_n|
 \le d_n\log\frac{CR_{n,\alpha}}\zeta
 =k_n\log L_n+O(d_n).
\]
Since \(e^\alpha t>1\), the union bound proves
\eqref{eq:slow-outside-shell}.
\end{proof}

The population gap in \eqref{eq:population-field} controls the entire inner
half-ball once the
microscopic curvature estimate is combined with a relative empirical-process
bound.
\begin{lemma}
\label{lem:slow-bulk}
Suppose \(3\le d_n\le(\log L_n)^2\) and \(s_n=O(1)\).  Then
\begin{equation}
\label{eq:slow-bulk}
 \Pbb\left\{
  \sup_{0<|w|\le R_{n,\alpha}/2}B_n(w)>1
 \right\}\longrightarrow0.
\end{equation}
\end{lemma}

\begin{proof}
At the origin, \(B_n(0)=1\), \(\nabla B_n(0)=0\), and Gaussian
sample-covariance concentration gives
\[
 \nabla^2B_n(0)
 =\frac1n\sum_i\frac{\widetilde Z_i\widetilde Z_i^\top}{\gamma_n}-I
 \preceq-\frac{\eps_n}{2}I
\]
with probability tending to one.  Since
\(\max_i|\widetilde Z_i|=O_{\Pbb}(\sqrt{L_n})\), Taylor's theorem makes
the inequality strict on
\(0<|w|\le r_{0,n}:=\eps_n/(16L_n^{3/2})\).

For completeness, we give the empirical-process estimate used beyond this
microscopic ball.  Before centering, put
\[
 H_n(w)=\frac1n\sum_i
 \exp\left\{w\cdot Q_i-\frac{|w|^2}{2\gamma_n}\right\},
 \qquad \E H_n(w)=1,
\]
and \(g_n(r)=1-e^{-\eps_nr^2/2}\).  On
\(r_{0,n}\le|w|\le R_{n,\alpha}/2\),
\begin{equation}
\label{eq:slow-bulk-process}
 \sup_w\frac{|H_n(w)-1|}{g_n(|w|)}\longrightarrow_{\Pbb}0.
\end{equation}
To verify this, split the range into squared-radius annuli
\(\mathcal A_j=\{r_j\le |w|\le R_j\}\) with \(R_j^2\le2r_j^2\), and put
\(g_j=g_n(r_j)\).  There are \(O(\log L_n)\) annuli.  On an annulus with
outer radius \(R=R_j\), clip the lognormal summands at
\[
 B_R=\exp\{R^2/(2\gamma_n)+RT_n\},
 \qquad
 T_n^2=C\{d_n\log L_n+\log g_n(r_{0,n})^{-1}\}.
\]
We first control both discarded tails.  For a fixed \(w\), exponential
tilting gives
\[
 \sup_{|w|\le R}
 \E\!\left[
  e^{w\cdot Q-|w|^2/(2\gamma_n)}
  \mathbf1_{\{e^{w\cdot Q-|w|^2/(2\gamma_n)}>B_R\}}
 \right]
 \le e^{-cT_n^2}.
\]
Indeed, tilt \(Q\) from
\(\mathcal N(0,\gamma_n^{-1}I)\) to
\(\mathcal N(w/\gamma_n,\gamma_n^{-1}I)\); the threshold is then at least
\(\sqrt{\gamma_n}T_n\) standard deviations in the \(w\)-direction.  For the
empirical envelope, use
\[
 \sup_{|w|\le R}
 e^{w\cdot q-|w|^2/(2\gamma_n)}
 =
 \begin{cases}
  e^{\gamma_n|q|^2/2},&\gamma_n|q|\le R,\\
  e^{R|q|-R^2/(2\gamma_n)},&\gamma_n|q|>R.
 \end{cases}
\]
Thus exceeding \(B_R\) forces \(|q|>R/\gamma_n+T_n\).  Completing the
square against the \(\mathcal N(0,\gamma_n^{-1}I)\) density gives
\[
 \E\left[
  \sup_{|w|\le R}
  e^{w\cdot Q-|w|^2/(2\gamma_n)}
  \mathbf1_{\{\sup_{|v|\le R}
    e^{v\cdot Q-|v|^2/(2\gamma_n)}>B_R\}}
 \right]
 \le
 C e^{C d_n\log L_n}
 \int_{T_n}^{\infty}e^{-cu^2}\,du.
\]
Taking the constant in \(T_n^2\) sufficiently large makes this
\(o(g_j/\log L_n)\).  Markov's inequality and a union bound over the
annuli control the empirical discarded tails at the same order.

It remains to control the clipped fields.  Fix \(\rho>0\).  On the event
\(\max_i|Q_i|\le C\sqrt{L_n}\), each is
\(C\sqrt{L_n}B_R\)-Lipschitz.  Use a net on \(\mathcal A_j\) of mesh
\[
 \frac{\rho g_j}{C\sqrt{L_n}B_R}.
\]
Its logarithmic cardinality is \(O(d_nL_n)\).  Because
\(R^2\le L_n/2+O(1)\), one has \(B_R\le n^{1/4+o(1)}\), while Bernstein's
inequality at deviation \(\rho g_j/4\) has exponent at least
\[
 \frac{c n g_n(r)^2}{e^{R^2/\gamma_n}+B_Rg_n(r)}
 \ge n^{1/2-o(1)}g_n(r_{0,n})^2.
\]
This dominates \(d_nL_n\).  The union bound over all net points and all
annuli shows that the probability that the left side of
\eqref{eq:slow-bulk-process} exceeds \(\rho\) tends to zero.  Since
\(\rho>0\) was arbitrary, \eqref{eq:slow-bulk-process} follows.

Now
\[
 B_n(w)
 =e^{-w\cdot\bar Z_n/\sqrt{\gamma_n}}
   e^{-\eps_n|w|^2/2}H_n(w).
\]
Moreover,
\[
 \sup_{r_{0,n}\le r\le R_{n,\alpha}/2}
 \frac{r|\bar Z_n|/\sqrt{\gamma_n}}{g_n(r)}
 \le
 C\frac{|\bar Z_n|}{\eps_nr_{0,n}}
 =o_{\Pbb}(1),
\]
because \(d_n\le(\log L_n)^2\),
\(|\bar Z_n|=O_{\Pbb}(\sqrt{d_n/n})\), and
\(\eps_n\ge c\log L_n/L_n\).  Thus
\eqref{eq:slow-bulk-process} preserves a fixed fraction of the population
gap \(g_n(|w|)\), proving \eqref{eq:slow-bulk}.
\end{proof}


\subsubsection{Moderately growing dimension}
\label{sec:moderate-dimensional-theorem}

We now prove \eqref{eq:outer-low-target} and \eqref{eq:inner-target} under
\eqref{eq:direct-truncation-range}. The
argument deliberately uses a deeper truncation than
\eqref{eq:A-conditions}: this makes every remaining low-cloud summand so
small that elementary Bernstein bounds survive a Euclidean net.

Put
\[
  \mathfrak h_n=\log\frac{d_nL_n}{\delta_n\eps_n}.
\]
The expansions in Section~\ref{sec:regime-map} and
\eqref{eq:direct-truncation-range} imply
\begin{equation}
\label{eq:moderate-scale-relations}
  \mathfrak h_n\asymp\log L_n,
  \qquad
  \frac{\eps_nL_n}{\mathfrak h_n}\longrightarrow\infty,
  \qquad
  \log\frac{L_n^3}{\eps_n^3}=O(\log L_n).
\end{equation}
For \(d_n=o(L_n)\), these statements follow from
\(\delta_n=1+o(1)\) and
\(\eps_nL_n\sim(d_n/2)\log(2L_n/d_n)\); in particular,
\[
 \mathfrak h_n=2\log L_n-\log\log(2L_n/d_n)+O(1).
\]
When \(d_n\asymp L_n\), both \(\eps_n\) and \(\delta_n\) stay bounded away
from zero and infinity.  When \(d_n/L_n\to\infty\),
\(\eps_n\to1\) and \(\delta_n\asymp\sqrt{L_n/d_n}\).  The last two
conditions in \eqref{eq:direct-truncation-range} then give all three relations in
\eqref{eq:moderate-scale-relations}.
They also permit deterministic sequences \(\omega_n\to\infty\) and
\begin{equation}
\label{eq:moderate-cutoff}
  A_n=\frac{\omega_n\mathfrak h_n}{\delta_n}
  =o(L_n\wedge d_n).
\end{equation}
Indeed, this is immediate when \(d_n=O(L_n)\); when \(d_n/L_n\to\infty\),
it is equivalent, up to logarithmic factors, to
\(d_n(\log d_n)^2=o(L_n^3)\). Finally set
\begin{equation}
\label{eq:moderate-inner-radius}
  r_n^2=\frac{40\mathfrak h_n}{\eps_n}=o(L_n).
\end{equation}

We shall also use a uniform count for the upper radial cloud.  It is a
moderate-deviation extension of Lemma~\ref{lem:gamma-hazard}.

\begin{lemma}
\label{lem:moderate-upper-count}
Under \eqref{eq:direct-truncation-range}, if \(A_n\to\infty\) and
\(A_n=o(L_n\wedge d_n)\), then
\[
 n\Pbb\{Y_n\ge b_n-a_nA_n\}
 =\exp\{(1+o(1))A_n\}.
\]
Consequently, the number of indices with
\((|Z_i|^2/2-b_n)/a_n\ge-A_n\) is at most \(e^{2A_n}\) with probability
tending to one.
\end{lemma}

\begin{proof}
On \([b_n-a_nA_n,b_n]\), the reciprocal Gamma hazard equals
\(a_n\{1+o(1)\}\) uniformly.  Indeed,
\eqref{eq:tail-density-ratio} gives this once
\(a_nA_n=o(b_n-k_n+1)\).  For \(k_n=O(L_n)\), that relation follows from
\(A_n=o(d_n\wedge L_n)\); for \(k_n/L_n\to\infty\), it follows from
\(a_n\asymp\sqrt{k_n/L_n}\) and \(A_n=o(L_n)\).  Integrating the
hazard gives the displayed tail ratio.  The count is binomial with mean
\(e^{(1+o(1))A_n}\), so Markov's inequality proves the second claim.
\end{proof}

The next proposition proves the low-cloud estimate with more precision than
is needed for the Gumbel law. The stronger factor \(L_n^{-10}\) leaves enough
slack to differentiate the field up to three times in the support-count
argument.

\begin{proposition}
\label{prop:moderate-low-cloud}
Suppose \eqref{eq:direct-truncation-range} holds and \(s_n=O(1)\). Choose \(A_n,r_n\)
as in \eqref{eq:moderate-cutoff}--\eqref{eq:moderate-inner-radius}. Then
\begin{equation}
\label{eq:moderate-low-cloud}
  \sup_{|w|\ge r_n}B_n^{\rm lo}(w)
  =O_{\Pbb}(\delta_nL_n^{-10}).
\end{equation}
The same conclusion holds after differentiating the bump field up to three
times, with the right side replaced by \(o_{\Pbb}(\delta_nL_n^{-7})\).
\end{proposition}

\begin{proof}
We first work before empirical centering. Put
\(\xi_{n,i}^{\circ}=(|Z_i|^2/2-b_n)/a_n\) and define
\[
 S_n^{\rm lo}(w)
 =\frac1n\sum_{i=1}^n
   \exp\left\{w\cdot\frac{Z_i}{\sqrt{\gamma_n}}-\frac{|w|^2}{2}\right\}
   \mathbf1_{\{\xi_{n,i}^{\circ}<-A_n+1\}}.
\]
Completing the square and using
\(\gamma_nL_n=b_n+a_ns_n\) show that every summand is bounded by
\begin{equation}
\label{eq:moderate-jump-bound}
  \mathfrak j_n
  =\exp\{-\delta_n(A_n-1+s_n)\}
  =\exp\{-(1+o(1))\omega_n\mathfrak h_n\}.
\end{equation}
For fixed \(w\),
\begin{equation}
\label{eq:moderate-low-mean}
  \E S_n^{\rm lo}(w)
  \le \exp\{-\eps_n|w|^2/2\}.
\end{equation}
Since a nonnegative summand bounded by \(\mathfrak j_n\) has variance at most
\(\mathfrak j_n\) times its mean, Bernstein's inequality gives
\begin{equation}
\label{eq:moderate-bernstein}
 \Pbb\{S_n^{\rm lo}(w)-\E S_n^{\rm lo}(w)>t\}
 \le
 \exp\left\{-c\frac{t^2}
   {\mathfrak j_n(\E S_n^{\rm lo}(w)+t)}\right\}.
\end{equation}

With probability tending to one, all score centers have norm at most
\(C\sqrt{L_n}\). Every local maximum of the low bump field lies in their
convex hull, so it is enough to consider \(|w|\le C\sqrt{L_n}\). On this
ball the field is multiplicatively \(C\sqrt{L_n}\)-Lipschitz:
\begin{equation}
\label{eq:multiplicative-lipschitz}
  S_n^{\rm lo}(w')
  \ge e^{-C\sqrt{L_n}|w-w'|}S_n^{\rm lo}(w).
\end{equation}
Take a net of mesh \(c/\sqrt{L_n}\); its logarithmic cardinality is at most
\(C d_n\log L_n\). For \(|w|\ge r_n\),
\eqref{eq:moderate-low-mean} is at most \(e^{-20\mathfrak h_n}\). Apply
\eqref{eq:moderate-bernstein} with
\(t=\delta_nL_n^{-10}\). Equations
\eqref{eq:moderate-scale-relations} and \eqref{eq:moderate-jump-bound} give
\[
  \frac{t}{\mathfrak j_n}\gg d_n\log L_n,
  \qquad
  \frac{t^2}{\mathfrak j_ne^{-20\mathfrak h_n}}\gg d_n\log L_n.
\]
The union bound and \eqref{eq:multiplicative-lipschitz} prove the uncentered
version of \eqref{eq:moderate-low-cloud}.

The proof of Lemma~\ref{lem:centered-radial-max} gives
\[
 \max_i|\xi_{n,i}-\xi_{n,i}^{\circ}|=o_{\Pbb}(1).
\]
Moreover,
\[
 B_n(w)
 =\exp\left\{-w\cdot\frac{\bar Z_n}{\sqrt{\gamma_n}}\right\}
   B_n^{\circ}(w),
 \qquad
 \sup_{|w|\le C\sqrt{L_n}}
 \left|w\cdot\frac{\bar Z_n}{\sqrt{\gamma_n}}\right|
 =O_{\Pbb}(L_n/\sqrt n)=o_{\Pbb}(\delta_n).
\]
Replacing \(A_n-1\) by \(A_n\) therefore proves the centered statement.
For a multi-index \(\beta\) of order \(j\le3\), differentiation gives
\[
  \left|\partial_w^\beta e^{-|w-q_i|^2/2}\right|
  \le C_j(1+|w-q_i|^j)e^{-|w-q_i|^2/2}.
\]
On a fixed enlargement of the convex hull, both \(|w|\) and
\(\max_i|q_i|\) are \(O_{\Pbb}(\sqrt{L_n})\). Repeating the net-and-Bernstein
argument used for the order-zero field with this polynomial envelope costs at most
\(L_n^{j/2}\); the base estimate
\(O_{\Pbb}(\delta_nL_n^{-10})\) is therefore
\(o_{\Pbb}(\delta_nL_n^{-7})\) for every \(j\le3\). Outside that enlargement,
the polynomially weighted Gaussian envelope is radially decreasing toward the convex hull
once the enlargement radius exceeds \(\sqrt3\). Its supremum is consequently
attained on the controlled boundary, proving the derivative estimate globally.
\end{proof}

The inner field requires a different argument because it equals one at the
origin. Sample-covariance curvature handles a microscopic ball, and the
bounded-jump estimate \eqref{eq:moderate-bernstein} handles the remaining
annuli.

\begin{proposition}
\label{prop:moderate-inner-field}
Under the assumptions of Proposition~\ref{prop:moderate-low-cloud},
\begin{equation}
\label{eq:moderate-inner-field}
  \Pbb\left\{\sup_{0<|w|\le r_n}B_n(w)<1\right\}
  \longrightarrow1.
\end{equation}
\end{proposition}

\begin{proof}
At the origin,
\[
  \nabla B_n(0)=0,
  \qquad
  \nabla^2B_n(0)=\frac1n\sum_iq_iq_i^\top-I.
\]
The Gaussian sample-covariance bound and
\eqref{eq:direct-truncation-range} imply
\begin{equation}
\label{eq:moderate-origin-curvature}
  \left\|\frac1n\sum_iq_iq_i^\top-\frac1{\gamma_n}I\right\|_{\rm op}
  =o_{\Pbb}(\eps_n),
  \qquad
  \nabla^2B_n(0)\preceq-\frac{\eps_n}{2}I
\end{equation}
with probability tending to one. On the same event
\(\max_i|q_i|=O(\sqrt{L_n})\), so the Hessian is
\(O(L_n^{3/2})\)-Lipschitz near zero. It follows that
\begin{equation}
\label{eq:moderate-microscopic-ball}
  B_n(w)<1
  \quad\text{for}\quad
  0<|w|\le r_{0,n}:=\frac{\eps_n}{16L_n^{3/2}}.
\end{equation}

It remains to consider \(r_{0,n}\le|w|\le r_n\). Put
\(g(r)=1-e^{-\eps_nr^2/2}\). In this range
\[
  g(r)\ge c\eps_nr_{0,n}^2
  \ge c\frac{\eps_n^3}{L_n^3}=:g_{0,n}.
\]
The scale relations above give
\(\log g_{0,n}^{-1}=O(\log L_n)\). Cover each dyadic radial annulus by a
net of mesh \(c g(r)/\sqrt{L_n}\). The total logarithmic cardinality is
\(O(d_n\log L_n)\), and
\eqref{eq:multiplicative-lipschitz} shows that a crossing of level one would
produce a net point at level at least \(1-g(r)/8\).

For the low cloud, apply \eqref{eq:moderate-bernstein} with
\(t=g(r)/4\). Uniformly over all annuli, its exponent is bounded below by
\[
  c\frac{g_{0,n}^2}{\mathfrak j_n}\gg d_n\log L_n.
\]
Its mean is at most \(1-g(r)\) by
\eqref{eq:moderate-low-mean}. Lemma~\ref{lem:moderate-upper-count}, with a
one-unit buffer, shows that the upper cloud contains at most
\(e^{3A_n}\) points with probability tending to one. On
\(|w|\le r_n=o(\sqrt{L_n})\), each corresponding bump is at most
\(e^{-L_n+o(L_n)}\), so their total contribution is
\(e^{-L_n+o(L_n)}=o(g_{0,n})\). Finally, empirical centering changes the
logarithm of the field by at most
\[
  |w|\frac{|\bar Z_n|}{\sqrt{\gamma_n}},
  \qquad
  \sup_{r\ge r_{0,n}}
  \frac{|\bar Z_n|/\sqrt{\gamma_n}}{\eps_nr}
  =o_{\Pbb}(1),
\]
where the last relation follows from
\(|\bar Z_n|/\sqrt{\gamma_n}=O_{\Pbb}(\sqrt{L_n/n})\) and
\[
 \frac{L_n^2}{\sqrt n\,\eps_n^2}=o(1),
\]
which follows from \eqref{eq:direct-truncation-range}. The union bound now excludes every crossing outside
the microscopic ball. Together with
\eqref{eq:moderate-microscopic-ball}, this proves
\eqref{eq:moderate-inner-field}.
\end{proof}

Together with the upper-cloud separation established at the start of this
subsection, Propositions~\ref{prop:moderate-low-cloud} and
\ref{prop:moderate-inner-field} exclude all subthreshold contacts in the
direct-truncation range.

\subsubsection{The logarithmic crossover}
\label{sec:constant-crossover}

The slowly growing and direct-truncation arguments cover the sparse-edge
regimes except the crossover scale
\(d_n\asymp L_n^3/(\log L_n)^2\); the larger entropy regime is handled
in Section~\ref{sec:simplex}. Here a direct entropy comparison is
unnecessarily delicate.  A spatial argument is more effective: distinct
sample centers are separated by order \(\sqrt{L_n}\), whereas a neighborhood
of radius \(\sqrt{\log L_n}\) already makes every nonnearest Gaussian bump
negligible.

\begin{lemma}
\label{lem:crossover-geometry}
Fix \(0<c_-<c_+<\infty\), put \(\ell_n=\log L_n\), and suppose
\[
 c_-\frac{L_n^3}{\ell_n^2}
 \le d_n\le
 c_+\frac{L_n^3}{\ell_n^2},
 \qquad s_n=O(1).
\]
Then
\begin{equation}
\label{eq:crossover-scales}
 \gamma_n\asymp\frac{L_n^2}{\ell_n^2},
 \qquad
 \delta_n:=\frac{a_n}{\gamma_n}\asymp\frac{\ell_n}{L_n},
 \qquad
 \eps_n\longrightarrow1.
\end{equation}
Moreover, with probability tending to one,
\begin{equation}
\label{eq:crossover-separation}
 \max_i\frac{|Z_i|}{\sqrt{\gamma_n}}\le C\sqrt{L_n},
 \qquad
 \min_{i\ne j}|q_i-q_j|^2\ge3L_n,
 \qquad
 \frac{|\bar Z_n|}{\sqrt{\gamma_n}}
 =O_{\Pbb}\!\left(\sqrt{\frac{L_n}{n}}\right).
\end{equation}
\end{lemma}

\begin{proof}
The scale relations follow from \eqref{eq:superlog-regime} and
\eqref{eq:epsilon-superlog}.  Also
\(|\bar Z_n|=O_{\Pbb}(\sqrt{d_n/n})\), which gives the last estimate in
\eqref{eq:crossover-separation}.  A Gamma Chernoff bound and a union bound
over \(i\) give the first estimate there.  Finally,
\[
 q_i-q_j=\frac{Z_i-Z_j}{\sqrt{\gamma_n}},
 \qquad
 \frac{|Z_i-Z_j|^2}{2}\sim\chi^2_{d_n}.
\]
Since \(d_n/\gamma_n\sim2L_n\), the lower-tail Chernoff bound for
\(\chi^2_{d_n}\) is at most \(e^{-c d_n}\) at the threshold
\(|q_i-q_j|^2=3L_n\).  This remains summable over the fewer than \(n^2\)
pairs because \(d_n/L_n\to\infty\).
\end{proof}

We next use this separation together with truncation of the independent,
uncentered field.  The truncation is only needed away from the sample
centers, where each individual bump is already small.

\begin{proposition}
\label{prop:crossover-exclusion}
Under the assumptions of Lemma~\ref{lem:crossover-geometry}, for
every fixed \(\eta>0\),
\[
 \Pbb\left\{
  \max_i\xi_{n,i}\le s_n-\eta,\quad
  \sup_{w\in\mathbb R^{d_n}}B_n(w)>1
 \right\}\longrightarrow0.
\]
\end{proposition}

\begin{proof}
Work on the high-probability events in Lemma~\ref{lem:crossover-geometry}
and on
\(\{\max_i\xi_{n,i}\le s_n-\eta\}\).  Then
\[
 \max_i\frac{|q_i|^2}{2}\le L_n-g_n,
 \qquad
 g_n=\eta\frac{a_n}{\gamma_n}\asymp_\eta\frac{\ell_n}{L_n}.
\]
 Put
 \[
  Q_i=\frac{Z_i}{\sqrt{\gamma_n}},
  \qquad c_n=\frac{\bar Z_n}{\sqrt{\gamma_n}},
  \qquad
  S_n(w)=\frac1n\sum_i e^{w\cdot Q_i-|w|^2/2}.
 \]
 Then \(q_i=Q_i-c_n\) and, exactly,
 \begin{equation}
 \label{eq:crossover-centering-identity}
  B_n(w)=e^{-w\cdot c_n}S_n(w),
  \qquad
  \E S_n(w)=e^{-\eps_n|w|^2/2}.
 \end{equation}
 Lemma~\ref{lem:crossover-geometry} gives
 \begin{equation}
 \label{eq:crossover-centering-small}
  \sup_{|w|\le\sqrt{2L_n}}|w\cdot c_n|
 =O_{\Pbb}(L_n/\sqrt n)=o_{\Pbb}(g_n).
\end{equation}
In particular,
\[
 |c_n|=O_{\Pbb}\!\left(\sqrt{L_n/n}\right)
 =o_{\Pbb}(L_n^{-3/2}),
\]
which will make centering negligible even in the smallest annulus below.

 We first exclude neighborhoods of the sample centers.  Write
 \(R_n=\sqrt{8\ell_n}\).  The bump representation and the radial buffer
 give
 \[
  h_i=\frac1n e^{|q_i|^2/2}\le e^{-g_n}\le1-\frac{g_n}{2}
 \]
 for all large \(n\).  If \(|w-q_i|\le R_n\), pairwise separation gives
 \[
  B_n(w)
  \le e^{-g_n}
  +n\exp\left\{-\frac12(\sqrt{3L_n}-R_n)^2\right\}
  =1-\frac{g_n}{2}+e^{-L_n/2+o(L_n)}<1.
 \]
 Thus a crossing cannot occur within distance \(R_n\) of any \(q_i\).

We next exclude a fixed ball around the origin.  Standard Gaussian
sample-covariance concentration and \(d_n/n\to0\) imply
\[
 \frac1n\sum_iq_iq_i^\top\preceq\frac14I_{d_n}
\]
with probability tending to one.  Hence
\(\nabla^2B_n(0)\preceq-3I_{d_n}/4\).  On the radial-buffer event,
\(\max_i|q_i|\le\sqrt{2L_n}\), so the operator norm of the third derivative
of \(B_n\) is at most \(CL_n^{3/2}\) throughout a ball of radius
\(c_0L_n^{-3/2}\), where \(c_0>0\) is a sufficiently small universal
constant.  Taylor's theorem therefore gives
\begin{equation}
\label{eq:crossover-microscopic}
 B_n(w)\le1-\frac{|w|^2}{4},
 \qquad 0<|w|\le r_{0,n}:=c_0L_n^{-3/2}.
\end{equation}

It remains to pass from this microscopic ball to a fixed radius.  On the
event \(\max_i|Q_i|\le C\sqrt{L_n}\),
\begin{equation}
\label{eq:crossover-log-lipschitz}
 \sup_{|w|\le3}|\nabla\log S_n(w)|\le C\sqrt{L_n}.
\end{equation}
Partition \(\{r_{0,n}\le|w|\le3\}\) into the \(O(\ell_n)\) annuli
\(\mathcal A_j=\{r_j\le|w|\le2r_j\}\), where \(r_j=2^jr_{0,n}\),
truncating the last annulus at radius \(3\).  For each \(j\), take a
deterministic net of \(\mathcal A_j\) with mesh
\[
 \rho_j=c_1\frac{r_j^2}{\sqrt{L_n}},
\]
where \(c_1>0\) is a sufficiently small constant.  The logarithm of the
total number of net points is \(O(d_n\ell_n)=O(L_n^3/\ell_n)\).

For a fixed net point \(v\), define the independent variables
\[
 Y_i(v)=\min\left\{
  \frac1n e^{v\cdot Q_i-|v|^2/2},
  U_n
 \right\},
 \qquad
 U_n=n^{-1}e^{C_1\sqrt{L_n}},
\]
with \(C_1\) chosen so that truncation is inactive on
\(\{\max_i|Q_i|\le C\sqrt{L_n}\}\) for \(|v|\le3\).  Since
\(\eps_n\ge1/2\) for all large \(n\),
\[
 \sum_i\E Y_i(v)
 \le e^{-\eps_n|v|^2/2}
 \le1-c_2|v|^2,
 \qquad |v|\le3,
\]
for a numerical \(c_2>0\), while
\(\sum_i\operatorname{Var}Y_i(v)\le U_n\).  Bernstein's inequality gives
\[
 \Pbb\left\{\sum_iY_i(v)>
  e^{-\eps_n|v|^2/2}+c_3r_j^2\right\}
 \le\exp\{-c r_j^4/U_n\}.
\]
The exponent is at least \(n^{1-o(1)}L_n^{-6}\), uniformly in \(j\), and
therefore dominates the total net entropy.  Thus, with probability tending
to one, \(S_n(v)\le1-c_4r_j^2\) at every such net point.

If \(B_n(w)\ge1\) at some \(w\in\mathcal A_j\), then
\eqref{eq:crossover-centering-identity} and
\(|c_n|=o_{\Pbb}(r_{0,n})\) give \(S_n(w)\ge1-o(r_j^2)\).
At the nearest net point \(v\), \eqref{eq:crossover-log-lipschitz} and the
choice of \(c_1\) then give \(S_n(v)>1-c_4r_j^2\), a contradiction.  In
combination with \eqref{eq:crossover-microscopic}, this proves
\begin{equation}
\label{eq:crossover-inner}
 \sup_{0<|w|\le3}B_n(w)<1
\end{equation}
with probability tending to one on the radial-buffer event.

We finally exclude the part of the field outside the fixed ball and the
center neighborhoods.  Let \(\mathcal G_n\) be a deterministic
\(1/4\)-net of \(B(0,\sqrt{2L_n})\).  Then
\[
 \log|\mathcal G_n|
 \le d_n\log(C\sqrt{L_n})
 =O(L_n^3/\ell_n).
\]
For \(v\in\mathcal G_n\) with \(|v|\ge11/4\), set
\[
 \widetilde Y_i(v)
 =
 \min\left\{
  \frac1n e^{v\cdot Q_i-|v|^2/2},
  \mathfrak b_n
 \right\},
 \qquad \mathfrak b_n=2L_n^{-15/4}.
\]
These variables are independent, and
\[
 \sum_i\E\widetilde Y_i(v)
 \le e^{-\eps_n|v|^2/2}<0.05
\]
for all large \(n\).  Since
\(\sum_i\operatorname{Var}\widetilde Y_i(v)\le0.05\mathfrak b_n\), Bernstein's
inequality and a union bound give
\begin{equation}
\label{eq:crossover-outer-net}
 \Pbb\left\{
  \max_{\substack{v\in\mathcal G_n\\|v|\ge11/4}}
  \sum_i\widetilde Y_i(v)>0.9
 \right\}
 \le
 \exp\{O(L_n^3/\ell_n)-cL_n^{15/4}\}
 =o(1).
\end{equation}

Suppose now that \(\sup_wB_n(w)>1\), and let \(w_*\) be a global
maximizer.  Lemma~\ref{lem:slow-truncated-shell} places \(w_*\) in the
convex hull of the \(q_i\)'s, hence in \(B(0,\sqrt{2L_n})\).  The preceding
two parts of the proof imply
\[
 |w_*|>3,\qquad \min_i|w_*-q_i|>R_n.
\]
Let \(v\in\mathcal G_n\) satisfy \(|v-w_*|\le1/4\).  The semiconvex
estimate \eqref{eq:semiconvex-net-transfer} gives
\[
 B_n(v)\ge e^{-1/32}B_n(w_*)>e^{-1/32}.
\]
Moreover, \(|v|\ge11/4\) and
\(\min_i|v-q_i|\ge R_n-1/4\).  The exact identity
\[
 \frac1n e^{v\cdot Q_i-|v|^2/2}
 =e^{v\cdot c_n}h_i e^{-|v-q_i|^2/2},
\]
together with the radial-buffer bound on \(h_i\) and
\eqref{eq:crossover-centering-small}, shows that every uncentered summand
at \(v\) is at most
\[
 e^{|v\cdot c_n|}e^{-(R_n-1/4)^2/2}
 \le \mathfrak b_n.
\]
Thus truncation is inactive and, by
\eqref{eq:crossover-centering-identity},
\[
 \sum_i\widetilde Y_i(v)=S_n(v)
 =e^{v\cdot c_n}B_n(v)>0.9
\]
for all large \(n\), contradicting \eqref{eq:crossover-outer-net}.
\end{proof}

The exclusion propositions in this section cover, respectively, slowly
growing dimensions, the direct truncation range, and the remaining
crossover band. They are combined with the entropy estimate in
Proposition~\ref{prop:buffered-support-below-simplex}; no separate
regime-by-regime limit theorem is needed.
